\originalchapter{核力与氘核}\label{ch:deuteron}
\originalsection{从强相互作用到有效核力}
质子之间有库仑排斥, 原子核却能束缚, 因而必须存在更强的短程吸引. 核子间有效相互作用不是基本夸克之间的颜色力的简单静电版本. 核子是颜色单态复合物, 低能核力是强相互作用的剩余效应. 它近似电荷无关, 有自旋与同位旋依赖, 极短距离存在排斥, 且带张量项. 说``电荷无关''指强核力部分的近似, 不包含质子间库仑力, 也不否认同位旋破缺.

氘核只含一个质子和一个中子, 是检验模型的最简单束缚核. 其结合能约 $2.225\,\mathrm{MeV}$, 总自旋宇称为 $1^+$. 磁矩和四极矩说明它并非纯粹的 $s$ 波中心势束缚态, 张量力使 ${}^3S_1$ 与 ${}^3D_1$ 混合. 下文的球方阱只给主导 $s$ 波近似, 不预测四极矩.

\originalsection{质心分离}
对两粒子 $H=p_1^2/(2m_1)+p_2^2/(2m_2)+V(|\mathbf r_1-\mathbf r_2|)$, 令
$M=m_1+m_2$, $\mu=m_1m_2/M$, $\mathbf R=(m_1\mathbf r_1+m_2\mathbf r_2)/M$, $\mathbf r=\mathbf r_1-\mathbf r_2$. 对应共轭动量为 $\mathbf P=\mathbf p_1+\mathbf p_2$ 与 $\mathbf p=(m_2\mathbf p_1-m_1\mathbf p_2)/M$. 反解为 $\mathbf p_1=(m_1/M)\mathbf P+\mathbf p$, $\mathbf p_2=(m_2/M)\mathbf P-\mathbf p$. 代入动能时混合项相消, 得
\begin{equation}
 H=\frac{P^2}{2M}+\left[\frac{p^2}{2\mu}+V(r)\right].
\end{equation}
束缚能属于相对运动哈密顿量. 质心可以任意自由运动, 因而实验室总能量不是单独一个离散束缚能级. 对氘核应使用约化质量 $\mu\simeq m_N/2$, 不能使用整个氘核质量代替.

\originalsection{方阱束缚态及归一化}
取 $V=-V_0$ ($r<R$)、$V=0$ ($r>R$), $E=-B<0$. $s$ 波满足 $u(0)=0$ 和远处可积, 因而
\begin{equation}
 u(r)=\begin{cases}A\sin kr,&0<r<R,\\ C\e^{-\kappa r},&r>R,\end{cases}
 \quad k=\frac{\sqrt{2\mu(V_0-B)}}\hbar,\quad\kappa=\frac{\sqrt{2\mu B}}\hbar.
\end{equation}
连续性给 $C=A\sin(kR)\e^{\kappa R}$, 导数连续给本征值条件 $k\cot(kR)=-\kappa$. 归一化为
\begin{equation}\label{eq:deuteron-norm}
 A^{-2}=\frac R2-\frac{\sin2kR}{4k}+\frac{\sin^2kR}{2\kappa}.
\end{equation}
因此阱外概率 $P_{\rm out}=A^2\sin^2(kR)/(2\kappa)$. 浅束缚时 $\kappa$ 小, 尾部较长, 相当大的概率落在有效核力范围之外. 粒子间距的均方值 $\langle r^2\rangle=\int r^2|u|^2\dd r$ 与单个核子相对于质心的均方半径不同; 对等质量两粒子后者为前者的 $1/4$.

题目同时给定 $V_0,R,B$ 时必须核对它们是否满足本征值条件. 近似参数可能有小残差, 不应把连续性和导数连续都声称精确成立. 做数值复现时可以解一个参数给出一致模型, 或报告固定题设数值的残差并保留近似性质.

\originalsection{束缚阈值与激发态}
令 $B\to0^+$, $l=0$ 匹配条件变成 $kR=(n+1/2)\pi$. 第一 $s$ 态出现于 $kR=\pi/2$, 第二径向 $s$ 态出现于 $3\pi/2$. 非零 $l$ 的离心项提高束缚阈值. 对方阱的 $l=1$ 首个阈值, 零能外部径向解 $u\propto r^{-1}$, 匹配 $u'/u=-1/R$ 得内球 Bessel 条件 $j_0(kR)=0$, 即 $kR=\pi$.

因此在相同阱深下, $\pi<\sqrt{2\mu V_0}R/\hbar<3\pi/2$ 可有 $p$ 态而无第二径向 $s$ 态. 实际氘核的相互作用还带自旋、同位旋与张量结构; 方阱的阈值比较解释其浅束缚性质, 不构成所有核力模型均无激发束缚态的一般定理.

双中子的低能 ${}^1S_0$ 通道在实验上没有束缚态. 同类费米子的交换限制迫使它进入单重态, 该通道吸引较弱; 存在大的负散射长度和接近阈值的虚态, 与``没有任何相互作用''不同. 核力强弱的实验输入是判断的必要部分.
